\documentclass[11pt,a4paper]{article}
\usepackage[T1]{fontenc}
\usepackage[utf8]{inputenc}
\usepackage[polish,provide=*]{babel}
\usepackage{lmodern}
\usepackage{amsmath,amssymb}
\usepackage[margin=2.5cm]{geometry}
\usepackage{enumitem}
\usepackage{booktabs,array}
\usepackage[version=4]{mhchem}
\usepackage{tikz}
\usepackage{tcolorbox}
\tcbuselibrary{breakable}

\newtcolorbox{wskazowki}{%
  breakable,
  colback=gray!6, colframe=gray!45, colbacktitle=gray!20, coltitle=black,
  fonttitle=\bfseries, title={Wskazówki}%
}

\newcounter{zadanie}
\newcommand{\zadanie}[1]{%
  \ifnum\value{zadanie}>0 \clearpage\fi
  \refstepcounter{zadanie}%
  \section*{Zadanie~\thezadanie.\ #1}%
}

\title{Termodynamika statystyczna: zadania przygotowawcze do egzaminu}
\author{}
\date{}

\begin{document}
\maketitle

\zadanie{Układ dwustanowy}

Rozważ układ $N$ nieoddziałujących cząstek o spinie $1/2$ w jednorodnym polu
magnetycznym $B$. Każda cząstka może zajmować jeden z dwóch poziomów
energetycznych:

\begin{center}
\begin{tikzpicture}
  \draw[->] (0,-0.3) -- (0,1.6) node[above] {energia cząstki};
  \draw[thick] (0.6,0) -- (3.6,0) node[right] {$-\varepsilon$};
  \draw[thick] (0.6,1) -- (3.6,1) node[right] {$+\varepsilon$};
\end{tikzpicture}
\end{center}

Układ jest izolowany. Jego stan makroskopowy określają całkowita energia $E$
i liczba cząstek $N$.

Wyprowadź następujące wielkości fizyczne:
\begin{enumerate}[label=(\alph*)]
  \item liczbę mikrostanów $\Omega(E,N)$,
  \item temperaturę $T$,
  \item energię jako funkcję temperatury, $E(T)$,
  \item entropię $S(E,N)$; naszkicuj wykres $S(E)$ i znajdź maksimum entropii,
  \item pojemność cieplną $C_V(T)$; naszkicuj jej wykres.
\end{enumerate}

\begin{wskazowki}
\begin{enumerate}[label=(\alph*)]
  \item \textit{Liczba mikrostanów.}
    \begin{itemize}
      \item Oznacz przez $N_+$ i $N_-$ liczby cząstek na poziomach $+\varepsilon$
        i $-\varepsilon$. Wyraź je przez $E$ i $N$.
      \item Na ile sposobów można wybrać $N_+$ cząstek spośród $N$?
    \end{itemize}
  \item \textit{Temperatura.}
    \begin{itemize}
      \item $1/T = (\partial S/\partial E)_N$, gdzie $S = k_B\ln\Omega$.
      \item Uprość $\ln\Omega$ wzorem Stirlinga, $\ln n! \approx n\ln n - n$.
      \item Wynik najprościej zapisać przez stosunek obsadzeń $N_-/N_+$.
    \end{itemize}
  \item \textit{Energia jako funkcja temperatury.}
    \begin{itemize}
      \item Z wyniku punktu~(b) wyznacz $N_+(T)$ i $N_-(T)$.
    \end{itemize}
  \item \textit{Entropia.}
    \begin{itemize}
      \item Oblicz $\partial S/\partial E$, żeby znaleźć maksimum entropii.
      \item Dla skrajnych energii $E = \pm N\varepsilon$ skorzystaj
        z $\lim_{n\to 0} n\ln n = 0$.
    \end{itemize}
  \item \textit{Pojemność cieplna.}
    \begin{itemize}
      \item Zróżniczkuj $E(T)$ z punktu~(c) przy stałym $\varepsilon$.
      \item Do szkicu wykresu: do czego dąży $C_V$ dla $k_BT \ll \varepsilon$
        i dla $k_BT \gg \varepsilon$?
    \end{itemize}
\end{enumerate}
\end{wskazowki}

\zadanie{Układ o równoodległych poziomach energetycznych}

Rozważamy opisany rozkładem kanonicznym układ znajdujący się w równowadze ze
zbiornikiem ciepła o temperaturze $T$. Poziomy energetyczne układu dane są przez
\begin{equation*}
  \varepsilon_j = j\,\varepsilon, \qquad j = 0, 1, 2, 3, \ldots, \infty ,
\end{equation*}
gdzie $\varepsilon > 0$ jest odległością między sąsiednimi poziomami
energetycznymi. Wiemy, że średnia energia tego układu to
\begin{equation*}
  \langle E \rangle = E(T) = \alpha\varepsilon ,
\end{equation*}
gdzie $\alpha > 0$ jest pewnym wyrażeniem.

\begin{enumerate}[label=(\alph*)]
  \item Pokaż, że temperatura $T$ rozważanego układu spełnia równanie
    \begin{equation*}
      \frac{1}{k_B T} = \frac{1}{\varepsilon}\ln\left(1 + \frac{1}{\alpha}\right).
    \end{equation*}
  \item Pokaż, że entropia układu dana jest wyrażeniem
    \begin{equation*}
      S = k_B(1 + \alpha)\ln(1 + \alpha) - k_B\,\alpha\ln\alpha .
    \end{equation*}
\end{enumerate}

\begin{wskazowki}
\begin{enumerate}[label=(\alph*)]
  \item \textit{Temperatura.}
    \begin{itemize}
      \item Zapisz kanoniczną funkcję podziału układu jako sumę po poziomach
        energetycznych:
        \begin{equation*}
          Z = \sum_{j=0}^{\infty} e^{-\varepsilon_j/k_BT} .
        \end{equation*}
      \item Kolejne wyrazy tej sumy tworzą szereg geometryczny. Jaki jest jego
        iloraz? Skorzystaj ze wzoru
        $\sum_{j=0}^{\infty} q^j = 1/(1 - q)$ dla $|q| < 1$.
      \item Korzystając z $Z$, oblicz średnią energię:
        \begin{equation*}
          E(T,V,N) = k_BT^2\,\frac{\partial}{\partial T}\ln Z(T,V,N) .
        \end{equation*}
    \end{itemize}
  \item \textit{Entropia.}
    \begin{itemize}
      \item Skorzystaj ze wzoru $S = k_B\ln Z + E/T$.
      \item Dzięki wynikowi punktu~(a) możesz wyrazić $\ln Z$ i $E/T$
        przez $\alpha$ zamiast przez $T$.
    \end{itemize}
\end{enumerate}
\end{wskazowki}

\zadanie{Obsadzenia poziomów rotacyjnych HCl}

W jakiej temperaturze liczba cząsteczek HCl w pierwszym rotacyjnym poziomie
wzbudzonym jest równa $1/e$ liczby cząsteczek w rotacyjnym poziomie
podstawowym? Energia poziomu rotacyjnego o liczbie kwantowej $J$ wynosi
\begin{equation*}
  \varepsilon_J = B\,J(J+1), \qquad J = 0, 1, 2, \ldots,
\end{equation*}
gdzie $B = \hbar^2/(2I)$ jest stałą rotacyjną, a $I$ jest momentem
bezwładności cząsteczki. Charakterystyczna temperatura rotacji HCl jest równa
\begin{equation*}
  \theta^\text{rot} = \frac{B}{k_B} = \frac{\hbar^2}{2 I k_B} = 15.02~\text{K} .
\end{equation*}
Prawidłowa odpowiedź: $T = 14.3$~K.

\begin{wskazowki}
\begin{itemize}
  \item Jaka jest degeneracja $g_J$ poziomu $J$? Ile wartości przyjmuje
    liczba kwantowa $M$ dla danego $J$?
  \item W rozkładzie kanonicznym liczba cząsteczek na poziomie $J$ jest
    proporcjonalna do $g_J\,e^{-\varepsilon_J/k_BT}$.
\end{itemize}
\end{wskazowki}

\zadanie{Jądro azotu w polu magnetycznym}

Jądro atomu azotu jest obdarzone momentem magnetycznym, który może przyjmować
trzy orientacje w polu magnetycznym $B$. Stanom tym odpowiadają trzy poziomy
energetyczne o energiach, kolejno,
\begin{equation*}
  \varepsilon_0 = -\gamma_N\hbar B, \qquad
  \varepsilon_1 = 0, \qquad
  \varepsilon_2 = +\gamma_N\hbar B .
\end{equation*}

\begin{enumerate}[label=(\alph*)]
  \item Wyprowadź funkcję podziału $Z(T)$ pojedynczego jądra azotu w polu
    magnetycznym.
  \item Wyprowadź średnią energię $E(T)$ pojedynczego jądra.
  \item Oblicz granice $\lim_{T\to 0} E(T)$ oraz $\lim_{T\to\infty} E(T)$.
  \item Wyprowadź wyrażenia na populacje (obsadzenia) poziomów wzbudzonych
    względem stanu podstawowego, zdefiniowane jako $f_1 = N_1/N_0$
    i $f_2 = N_2/N_0$, gdzie $N_0$ oznacza liczbę jąder w stanie podstawowym $\varepsilon_0$,
    $N_1$ jest liczbą jąder w pierwszym stanie wzbudzonym $\varepsilon_1$,
    a $N_2$ jest liczbą jąder w drugim stanie wzbudzonym $\varepsilon_2$.
  \item Oblicz $f_1$ i $f_2$ w temperaturach:
    (i) $T = 1.0$~K, (ii) $T = 298$~K, (iii) $T = \infty$.
\end{enumerate}

Stałe: $\gamma_N\hbar = 2.04\times 10^{-27}$~J/T, \quad $B = 20$~T, \quad
$k_B = 1.381\times 10^{-23}$~J/K.

\zadanie{Pojemność cieplna gazów w temperaturze pokojowej}

W tabeli podano geometrię oraz charakterystyczne temperatury rotacji
$\theta^\text{rot}$ i oscylacji $\theta^\text{vib}$ dla atomu neonu i wybranych
cząsteczek. Cząsteczki liniowe mają jedną temperaturę rotacji, a nieliniowe
trzy. Liczba w nawiasie przy $\theta^\text{vib}$ oznacza degenerację, czyli
liczbę drgań normalnych o tej samej temperaturze charakterystycznej.

\begin{center}
\small
\renewcommand{\arraystretch}{1.15}
\begin{tabular}{l l l >{\raggedright\arraybackslash}p{5.7cm}}
  \toprule
  Gaz & Geometria & $\theta^\text{rot}$ (K) & $\theta^\text{vib}$ (K) \\
  \midrule
  \ce{Ne}    & pojedynczy atom            & ---                       & --- \\
  \ce{CO2}   & liniowa                    & 0.561                     & 3360, 954~(2), 1890 \\
  \ce{H2O}   & nieliniowa (kątowa)        & 40.1, 20.9, 13.4          & 5360, 5160, 2290 \\
  \ce{NH3}   & nieliniowa (piramidalna)   & 13.6, 13.6, 8.92          & 4800, 1360, 4880~(2), 2330~(2) \\
  \ce{ClO2}  & nieliniowa (kątowa)        & 2.50, 0.478, 0.400        & 1360, 640, 1600 \\
  \ce{SO2}   & nieliniowa (kątowa)        & 2.92, 0.495, 0.422        & 1660, 750, 1960 \\
  \ce{N2O}   & liniowa                    & 0.603                     & 3200, 850~(2), 1840 \\
  \ce{CH4}   & nieliniowa (tetraedryczna) & 7.54, 7.54, 7.54          & 4170, 2180~(2), 4320~(3), 1870~(3) \\
  \ce{CH3Cl} & nieliniowa (tetraedryczna) & 7.32, 0.637, 0.637        & 4270, 1950, 1050, 4380~(2), 2140~(2), 1460~(2) \\
  \ce{CCl4}  & nieliniowa (tetraedryczna) & 0.0823, 0.0823, 0.0823    & 660, 310~(2), 1120~(3), 450~(3) \\
  \bottomrule
\end{tabular}
\end{center}

\begin{enumerate}[label=(\alph*)]
  \item Jaką pojemność cieplną $C_V/(Nk_B)$ przewidujesz dla wymienionych
    gazów w temperaturze pokojowej, $T = 298$~K? Każdy gaz traktuj jako
    doskonały. Przyjmij, że stopnie swobody o temperaturze charakterystycznej
    $\theta \ll T$ są w reżimie klasycznym, a stopnie swobody
    o $\theta \gg T$ są zamrożone.
  \item Dla których gazów oszacowanie z punktu~(a) jest najmniej pewne?
    Wskaż drgania, dla których $\theta^\text{vib}$ jest porównywalne z $T$.
\end{enumerate}

\begin{wskazowki}
    \begin{itemize}
      \item Pojemność cieplna gazu doskonałego jest sumą wkładów translacji,
        rotacji i drgań. Wzbudzenia elektronowe pomiń.
      \item Zasada ekwipartycji: każdy kwadratowy wyraz w klasycznej energii
        cząsteczki daje wkład $k_BT/2$ do jej średniej energii, czyli $k_B/2$
        do $C_V/N$.
      \item Translacje: ile kwadratowych wyrazów ma energia kinetyczna ruchu
        środka masy?
      \item Rotacje: porównaj $\theta^\text{rot}$ z $T$. Ile niezależnych
        rotacji ma cząsteczka liniowa, a ile nieliniowa? Pojedynczy atom nie
        wnosi wkładu rotacyjnego.
      \item Drgania: cząsteczka złożona z $n$ atomów ma $3n - 5$ (liniowa) lub
        $3n - 6$ (nieliniowa) drgań normalnych. Sprawdź tę liczbę w tabeli,
        uwzględniając degeneracje.
      \item Energia drgania w reżimie klasycznym ma dwa kwadratowe wyrazy:
        kinetyczny i potencjalny. Drganie zamrożone nie wnosi wkładu.
    \end{itemize}
\end{wskazowki}

\zadanie{Elektronowy wkład do entropii tlenu}

Podstawowy poziom energii elektronowej cząsteczki \ce{O2} ma degenerację
$g_0 = 3$. Pierwszy wzbudzony poziom energii elektronowej ma degenerację
$g_1 = 2$ i leży $\tilde\nu = 7824$~cm$^{-1}$ powyżej poziomu podstawowego.
Wyższe poziomy elektronowe pomiń. Tlen traktuj jako gaz doskonały.

\begin{enumerate}[label=(\alph*)]
  \item Oblicz elektronowy wkład do entropii molowej tlenu, $S^\text{el}_m$,
    w temperaturze $T = 5000$~K.
  \item Oblicz $S^\text{el}_m$ w tej samej temperaturze, uwzględniając tylko
    poziom podstawowy. Porównaj oba wyniki.
\end{enumerate}

{\raggedright
Stałe: \mbox{$h = 6.626\times 10^{-34}$~J\,s},
\mbox{$c = 2.998\times 10^{10}$~cm/s},
\mbox{$k_B = 1.381\times 10^{-23}$~J/K},
\mbox{$R = 8.314$~J/(mol\,K)}.

Prawidłowa odpowiedź: (a)~$S^\text{el}_m = 10.9$~J/(mol\,K),
(b)~$S^\text{el}_m = 9.1$~J/(mol\,K).\par}

\begin{wskazowki}
\begin{itemize}
  \item Zapisz elektronową funkcję podziału jednej cząsteczki, $q^\text{el}$.
    Przyjmij energię poziomu podstawowego równą zeru i pamiętaj
    o degeneracjach.
  \item Cząsteczki nie oddziałują, więc dla $N$ cząsteczek
    $Z^\text{el} = (q^\text{el})^N$ (czynnik $1/N!$ poznany na wykładzie
    uwzględnia się w części translacyjnej).
  \item Entropia: $S^\text{el} = k_B\ln Z^\text{el} + E^\text{el}(T)/T$, gdzie
    $E^\text{el}(T) = k_BT^2\,\partial\ln Z^\text{el}/\partial T$.
  \item Energia wzbudzenia: $\varepsilon = hc\tilde\nu$.
\end{itemize}
\end{wskazowki}


\end{document}
